\scanpage{138}\mspage{12}
% archive sheet number: 137
Ceci posé, pour toute flèche $f$ de $M$, et tout ordinal, on a
\[
i_\alpha(f)\in\Phi_0.
\]
\deleted{\uline{Corollaire 2} \inserted{Supposons les sources des $i\in C_0$ \unsure{soient} accessibles.} Dans les conditions du corollaire 1, le couple $(\Phi_0,\Psi)$ satisfait : $\Phi_0\longleftrightarrow\Psi$ et à la condition de factorisation de Quillen. \deleted{Si $\Phi_0$ est de plus stable par facteurs directs, \unsure{alors|donc} $\Phi_0=\Phi$.} \unsure{de plus stable par facteurs directs}
\margin{\deleted{Donc $\Phi_0:$ $\Phi_0\unsure{\subset|=}\Psi^*$; satisfait $\Phi_0\longleftrightarrow\Psi$}}
En effet, comme le couple $(\Phi_0,\Psi)$ \deleted{satisfait} \inserted{\unsure{est satisfaite}} la condition de factorisation de Quillen par le cor. 1 et le th. 1 \inserted{b)} (\unsure{comme pour a}), \unsure{il sera}}

\uline{Corollaire 2} Supposons que les sources des flèches $i\in C_0$ soient accessibles, et soit $\Phi_0$ comme dans le cor. 1 (\unsure{p.e.|p.a.} $\Phi_0=C'''_0$). Alors le couple $(\Phi_0,\Psi)$ \deleted{satisfait} est un \unsure{précouple|pré-couple} de Quillen, i.e. a) $\Phi_0\longleftrightarrow\Psi$ i.e. $\Psi\subset\Phi_{0*}$ \unsure{Ce|ce|(ce} qui revient au $\unsure{=|\simeq}$, $\Phi_0\subset\Psi^*$) et b) on a la propriété de factorisation de Quillen. Cela implique donc que $\Phi$ \inserted{$(=\widetilde C_0=\Phi_0)$} est \unsure{formé|fermé} des \deleted{\unsure{f}} flèches qui sont facteurs directs des flèches de $\Phi_0$. Donc si $\Phi_0$ est stable par facteurs directs, on a $\Phi_0=\Phi$, i.e. $\Phi_0$ est clos à gauche.

2. \uline{\unsure{Variation|Variations} des \deleted{\unsure{les|la}} factorisations} $f=p_\alpha i_\alpha$ \uline{avec} \uline{la partie} $C_0$ génératrice.

Je dis que le système inductif ordinal de foncteurs $\unsure{P_\alpha|p_\alpha}=P_\alpha^{C_0}:\uline{\mathrm{Fl}}\,M\mathrel{\text{---}}\uline{\mathrm{Fl}}\,M$,
